\shead{Venn diagrams with two circles}
% ---- chapter 12 helpers: Venn diagrams with two and three circles ----
\newcommand{\stpVennII}[6]{% #1 #2 circle labels; #3 only A, #4 both, #5 only B, #6 outside (text)
\begin{tikzpicture}[font=\small]
\fill[tmBlue!14] (0,0) circle (1.3);
\fill[tmOrange!18] (1.7,0) circle (1.3);
\begin{scope}\clip (0,0) circle (1.3);\fill[tmGreen!35] (1.7,0) circle (1.3);\end{scope}
\draw[line width=0.9pt,tmBlue!80!black] (0,0) circle (1.3);
\draw[line width=0.9pt,tmOrange!85!black] (1.7,0) circle (1.3);
\draw[line width=0.8pt,Slate,rounded corners=4pt] (-2.0,-2.0) rectangle (3.7,1.95);
\node[font=\footnotesize\bfseries,text=tmBlue!80!black] at (-0.55,1.62) {#1};
\node[font=\footnotesize\bfseries,text=tmOrange!85!black] at (2.25,1.62) {#2};
\node at (-0.45,0) {#3}; \node at (0.85,0) {#4}; \node at (2.15,0) {#5};
\node[anchor=south east,font=\footnotesize] at (3.6,-1.95) {#6};
\end{tikzpicture}}
\newcommand{\stpVennIII}[9]{% #1 #2 #3 labels A B C; #4 only A, #5 only B, #6 only C, #7 A&B only, #8 A&C only, #9 B&C only
\def\stpVa{#1}\def\stpVb{#2}\def\stpVc{#3}\def\stpVd{#4}\def\stpVe{#5}\def\stpVf{#6}\def\stpVg{#7}\def\stpVh{#8}\def\stpVi{#9}\stpVennIIIx}
\newcommand{\stpVennIIIx}[2]{% #1 all three, #2 outside (text)
\begin{tikzpicture}[font=\small]
\fill[tmBlue!30,fill opacity=0.55] (-0.7,0.4) circle (1.15);
\fill[tmOrange!35,fill opacity=0.55] (0.7,0.4) circle (1.15);
\fill[tmGreen!30,fill opacity=0.55] (0,-0.8) circle (1.15);
\draw[line width=0.9pt,tmBlue!80!black] (-0.7,0.4) circle (1.15);
\draw[line width=0.9pt,tmOrange!85!black] (0.7,0.4) circle (1.15);
\draw[line width=0.9pt,tmGreen!70!black] (0,-0.8) circle (1.15);
\draw[line width=0.8pt,Slate,rounded corners=4pt] (-4.0,-2.3) rectangle (4.0,1.85);
\node[font=\footnotesize\bfseries,text=tmBlue!80!black,anchor=east] at (-1.85,1.0) {\stpVa};
\node[font=\footnotesize\bfseries,text=tmOrange!85!black,anchor=west] at (1.85,1.0) {\stpVb};
\node[font=\footnotesize\bfseries,text=tmGreen!60!black,anchor=west] at (1.2,-1.65) {\stpVc};
\node at (-1.25,0.65) {\stpVd}; \node at (1.25,0.65) {\stpVe}; \node at (0,-1.45) {\stpVf};
\node at (0,0.98) {\stpVg}; \node at (-0.72,-0.45) {\stpVh}; \node at (0.72,-0.45) {\stpVi};
\node at (0,0.05) {#1};
\node[anchor=south west,font=\footnotesize] at (-3.9,-2.25) {#2};
\end{tikzpicture}}
\begin{knowbox}
A \textbf{Venn diagram} shows groups as closed curves (circles) inside a rectangle that stands for everybody.
\begin{center}
\stpVennII{Group A}{Group B}{\scriptsize only A}{\scriptsize both}{\scriptsize only B}{in neither group}
\end{center}
\begin{itemize}[leftmargin=1.4em,itemsep=3pt]
\item The \emph{overlap} of the two circles holds those who belong to \textbf{both} groups.
\item The rest of each circle holds those who belong to \textbf{only} that group.
\item Inside the rectangle but outside both circles are those who belong to \textbf{neither} group.
\end{itemize}
When we add A $+$ B, the members of both groups are counted \emph{twice}. So:
\begin{center}
\renewcommand{\arraystretch}{1.3}
\begin{tabular}{@{}r@{\ $=$\ }l@{}}
\textbf{at least one group} & A $+$ B $-$ both\\
\textbf{both} & A $+$ B $-$ (at least one group)\\
\textbf{only A} & A $-$ both\\
\textbf{neither} & everybody $-$ (at least one group)\\
\end{tabular}
\end{center}
\end{knowbox}
\shead{Venn diagrams with three circles}
\begin{knowbox}
Three circles A, B, C make $7$ regions inside them (plus the region outside). Fill in the diagram \textbf{from the centre outwards}:
\begin{enumerate}[leftmargin=1.6em,itemsep=2pt]
\item the centre: those in all three groups;
\item each overlap of two circles: (in both groups) $-$ (in all three);
\item each ``only'' region: (whole circle) $-$ (the three numbers already written in it).
\end{enumerate}
Adding all the regions gives the number in at least one group. In one line:
\begin{center}
at least one $=$ A $+$ B $+$ C $-$ (A and B) $-$ (B and C) $-$ (C and A) $+$ (all three).
\end{center}
\end{knowbox}
\shead{Working backwards with a table}
\begin{knowbox}
When things are passed from one person (or box) to another several times and we know the \textbf{final} result, we work \textbf{backwards} from the end:
\begin{itemize}[leftmargin=1.4em,itemsep=3pt]
\item Passing things around does \emph{not} change the total. If at the end everyone has the same, each has total $:$ number of people.
\item Undo the steps in reverse order: whoever \emph{gave} gets the amount back ($+$), whoever \emph{received} gives it back ($-$).
\item ``Gives as many as the other already has'' makes the receiver's amount \emph{double}; going back we \emph{halve} it.
\item A \textbf{table} with one row for each moment keeps the work tidy. Check by working forwards.
\end{itemize}
\end{knowbox}
\gap
\begin{exbox}{\ExWord\ 1}
A group of $30$ pupils each joins at least one of two clubs: $18$ join the Chess club and $19$ join the Drawing club. How many pupils join both clubs? How many join only one club?
\guide
Adding $18+19=37$ counts the pupils in both clubs twice, so
\par both clubs: $37-30=7$ pupils.\par
Only Chess: $18-7=11$; only Drawing: $19-7=12$; only one club: $11+12=23$ pupils.
\begin{center}
\stpVennII{Chess $18$}{Drawing $19$}{$11$}{$\mathbf{7}$}{$12$}{$30$ pupils}
\end{center}
\dapso{$7$ pupils; $23$ pupils.}
\end{exbox}
\gap
\begin{exbox}{\ExWord\ 2}
Class 5C has $36$ pupils. $20$ pupils can swim, $15$ can play badminton and $8$ can do both. How many pupils can do neither?
\guide
\emph{Method 1.} At least one: $20+15-8=27$ pupils.\par Neither: $36-27=9$ pupils.\par
\emph{Method 2.} Only swim: $20-8=12$; only badminton: $15-8=7$.\par At least one: $12+7+8=27$; neither: $36-27=9$.
\begin{center}
\stpVennII{Swim $20$}{Badminton $15$}{$12$}{$8$}{$7$}{neither: $\mathbf{9}$}
\end{center}
\dapso{$9$ pupils.}
\end{exbox}
\gap
\begin{exbox}{\ExWord\ 3}
Every pupil in a group likes at least one of football, swimming and chess. $25$ like football, $20$ like swimming and $18$ like chess; $8$ like football and swimming, $6$ like swimming and chess, $7$ like football and chess, and $3$ like all three.\par
\pa{a} How many like only football? \quad \pa{b} How many pupils are in the group?
\guide
Fill in from the centre: all three: $3$.\par
Football and swimming only: $8-3=5$; \ swimming and chess only: $6-3=3$; \ football and chess only: $7-3=4$.\par
\pa{a} Only football: $25-5-4-3=13$.\par
Only swimming: $20-5-3-3=9$; \ only chess: $18-4-3-3=8$.
\begin{center}
\stpVennIII{Football $25$}{Swimming $20$}{Chess $18$}{$\mathbf{13}$}{$9$}{$8$}{$5$}{$4$}{$3$}{$3$}{}
\end{center}
\pa{b} $13+9+8+5+4+3+3=45$ pupils.\par
(Check: $25+20+18-8-6-7+3=45$.)
\dapso{\pa{a} $13$ pupils. \pa{b} $45$ pupils.}
\end{exbox}
\gap
\begin{exbox}{\ExWord\ 4}
An, Bình and Chi have $75$ stickers altogether. An gives Bình $6$ stickers, then Bình gives Chi $9$ stickers, then Chi gives An $4$ stickers. Now the three have the same number. How many stickers did each have at first?
\guide
The total stays $75$, so at the end each has $75:3=25$ stickers. Work backwards, one row per step:
\begin{center}
\renewcommand{\arraystretch}{1.25}
\begin{tabular}{|l|c|c|c|}
\hline
 & An & Bình & Chi\\\hline
At the end & $25$ & $25$ & $25$\\\hline
Before ``Chi gives An $4$'' & $25-4=21$ & $25$ & $25+4=29$\\\hline
Before ``Bình gives Chi $9$'' & $21$ & $25+9=34$ & $29-9=20$\\\hline
Before ``An gives Bình $6$'' (at first) & $21+6=\mathbf{27}$ & $34-6=\mathbf{28}$ & $\mathbf{20}$\\\hline
\end{tabular}
\end{center}
Check: $27+28+20=75$.
\dapso{An $27$, Bình $28$, Chi $20$ stickers.}
\end{exbox}
\gap
\begin{exbox}{\ExWord\ 5}
Lan, Mai and Cúc have $120$ shells. First Lan gives Mai and Cúc each as many shells as each of them already has. Then Mai does the same for Lan and Cúc, and finally Cúc does the same for Lan and Mai. Now all three have the same number. How many shells did each have at first?
\guide
At the end each has $120:3=40$. Each time, the two receivers \emph{double}; going back we halve them, and the giver gets back what they received.
\begin{center}
\renewcommand{\arraystretch}{1.25}
\begin{tabular}{|l|c|c|c|}
\hline
 & Lan & Mai & Cúc\\\hline
At the end & $40$ & $40$ & $40$\\\hline
Before Cúc gives & $20$ & $20$ & $40+20+20=80$\\\hline
Before Mai gives & $10$ & $20+10+40=70$ & $40$\\\hline
At first (before Lan gives) & $10+35+20=\mathbf{65}$ & $\mathbf{35}$ & $\mathbf{20}$\\\hline
\end{tabular}
\end{center}
Check forwards: $(65,35,20)\to(10,70,40)\to(20,20,80)\to(40,40,40)$.
\dapso{Lan $65$, Mai $35$, Cúc $20$ shells.}
\end{exbox}
\dividerbar{Practice}
\begin{probbox}{1}{\relax}
Every pupil of class 5A got at least one mark $10$ in Maths or Vietnamese. $21$ pupils got a $10$ in Maths, $17$ in Vietnamese, and $6$ got a $10$ in both. How many pupils are there in class 5A?
\end{probbox}
\probsol{
$21+17-6=32$ pupils.
\dapso{$32$ pupils.}}
\gap
\begin{probbox}{2}{\relax}
Each of $40$ pupils likes reading or drawing (or both). $28$ like reading and $22$ like drawing. How many like both?
\end{probbox}
\probsol{
$28+22=50$ counts those who like both twice.\par Both: $50-40=10$ pupils.
\dapso{$10$ pupils.}}
\gap
\begin{probbox}{3}{\relax}
In a group of $45$ tourists, $30$ speak English, $20$ speak French and $5$ speak neither language.\par
\pa{a} How many speak both? \quad \pa{b} How many speak only English? \quad \pa{c} only French?
\end{probbox}
\probsol{
At least one language: $45-5=40$.\par
\pa{a} Both: $30+20-40=10$.\par
\pa{b} Only English: $30-10=20$.\par
\pa{c} Only French: $20-10=10$.
\dapso{\pa{a} $10$. \pa{b} $20$. \pa{c} $10$.}}
\gap
\needspace{7.5cm}
\begin{probbox}{4}{\relax}
The Venn diagram shows what the people at a party drink.
\begin{center}
\stpVennII{Tea}{Coffee}{$14$}{$6$}{$9$}{neither: $3$}
\end{center}
\pa{a} How many drink tea? \quad \pa{b} How many drink coffee? \quad \pa{c} How many people are at the party?
\end{probbox}
\probsol{
\pa{a} Tea: $14+6=20$.\par
\pa{b} Coffee: $6+9=15$.\par
\pa{c} $14+6+9+3=32$ people.
\dapso{\pa{a} $20$. \pa{b} $15$. \pa{c} $32$ people.}}
\gap
\needspace{7.5cm}
\begin{probbox}{5}{\relax}
Each member of a sports team plays football or basketball (or both), as the diagram shows.
\begin{center}
\stpVennII{Football}{Basketball}{$12$}{$4$}{$9$}{}
\end{center}
\pa{a} How many members does the team have? \quad \pa{b} How many play basketball?
\end{probbox}
\probsol{
\pa{a} $12+4+9=25$ members.\par
\pa{b} $4+9=13$ members.
\dapso{\pa{a} $25$. \pa{b} $13$.}}
\gap
\begin{probbox}{6}{\relax}
A maths team of $20$ pupils is given two hard problems; everybody solves at least one. $3$ pupils cannot solve problem 1 and $5$ cannot solve problem 2. How many solve both problems?
\end{probbox}
\probsol{
The $3$ who cannot solve problem 1 solve only problem 2.\par The $5$ who cannot solve problem 2 solve only problem 1.\par
Both: $20-3-5=12$ pupils.
\dapso{$12$ pupils.}}
\gap
\begin{probbox}{7}{\relax}
Class 5D has $38$ pupils. $15$ sign up for swimming, $20$ for football and $7$ for both. How many sign up for neither?
\end{probbox}
\probsol{
At least one: $15+20-7=28$.\par Neither: $38-28=10$ pupils.
\dapso{$10$ pupils.}}
\gap
\begin{probbox}{8}{\relax}
In a class of $35$ pupils, $14$ have a dog, $11$ have a cat and $5$ have both. How many pupils have neither a dog nor a cat?
\end{probbox}
\probsol{
At least one pet: $14+11-5=20$.\par Neither: $35-20=15$ pupils.
\dapso{$15$ pupils.}}
\gap
\begin{probbox}{9}{\relax}
In a village, $60$ families were asked: $42$ have a television, $35$ have a computer, and every family has at least one of them.\par
\pa{a} How many have both? \quad \pa{b} only a television? \quad \pa{c} only a computer?
\end{probbox}
\probsol{
\pa{a} $42+35-60=17$ families.\par
\pa{b} $42-17=25$ families.\par
\pa{c} $35-17=18$ families.
\dapso{\pa{a} $17$. \pa{b} $25$. \pa{c} $18$.}}
\gap
\begin{probbox}{10}{\relax}
How many of the numbers from $1$ to $30$ are divisible by $2$ or by $3$?
\end{probbox}
\probsol{
Divisible by $2$: $30:2=15$ numbers. Divisible by $3$: $30:3=10$ numbers.\par
Divisible by both, i.e.\ by $6$: $6,\ 12,\ 18,\ 24,\ 30$: $5$ numbers.
\begin{center}
\stpVennII{Divisible by $2$}{Divisible by $3$}{$10$}{$5$}{$5$}{}
\end{center}
By $2$ or by $3$: $15+10-5=20$ numbers.
\dapso{$20$ numbers.}}
\gap
\begin{probbox}{11}{\relax}
How many of the numbers from $1$ to $50$ are divisible neither by $2$ nor by $5$?
\end{probbox}
\probsol{
By $2$: $25$ numbers; by $5$: $10$ numbers; by both (by $10$): $5$ numbers.\par
By $2$ or $5$: $25+10-5=30$.\par Neither: $50-30=20$ numbers.
\dapso{$20$ numbers.}}
\gap
\begin{probbox}{12}{\relax}
In a music group, $16$ pupils play the piano, $10$ play the guitar, and $22$ pupils play at least one of them. How many play both? How many play only the piano?
\end{probbox}
\probsol{
Both: $16+10-22=4$ pupils.\par Only piano: $16-4=12$ pupils.
\dapso{$4$ pupils; $12$ pupils.}}
\gap
\begin{probbox}{13}{\relax}
$18$ pupils of a class entered the singing contest, $13$ entered the drawing contest, and $25$ pupils entered at least one contest. How many pupils entered exactly one contest?
\end{probbox}
\probsol{
Both contests: $18+13-25=6$ pupils.\par Exactly one: $25-6=19$ pupils.
\dapso{$19$ pupils.}}
\gap
\begin{probbox}{14}{\relax}
$100$ pupils were asked about breakfast: $58$ like phở, $47$ like bánh mì and $15$ like neither. How many like both?
\end{probbox}
\probsol{
At least one: $100-15=85$.\par Both: $58+47-85=20$ pupils.
\dapso{$20$ pupils.}}
\gap
\needspace{7.5cm}
\begin{probbox}{15}{\relax}
The diagram shows how many pupils of a class are in the Maths, English and Art clubs.
\begin{center}
\stpVennIII{Maths}{English}{Art}{$8$}{$7$}{$6$}{$3$}{$2$}{$4$}{$1$}{no club: $5$}
\end{center}
\pa{a} How many pupils are in the Maths club? \quad \pa{b} How many are in exactly one club?\par
\pa{c} How many pupils are there in the class?
\end{probbox}
\probsol{
\pa{a} Every number inside the Maths circle: $8+3+2+1=14$.\par
\pa{b} $8+7+6=21$.\par
\pa{c} $8+7+6+3+2+4+1+5=36$.
\dapso{\pa{a} $14$. \pa{b} $21$. \pa{c} $36$ pupils.}}
\gap
\begin{probbox}{16}{\relax}
Every pupil of a contest team takes at least one of Maths, English and Art: $25$ take Maths, $20$ take English and $15$ take Art. $6$ take Maths and English, $4$ take English and Art, $5$ take Maths and Art, and nobody takes all three. How many pupils are in the team?
\end{probbox}
\probsol{
Adding $25+20+15=60$ counts each pupil who takes two subjects twice.\par
Team: $60-6-4-5=45$ pupils.
\dapso{$45$ pupils.}}
\gap
\begin{probbox}{17}{\relax}
$40$ pupils were asked about fruit: $20$ like apples, $18$ bananas, $16$ oranges; $7$ like apples and bananas, $6$ bananas and oranges, $5$ apples and oranges, and $2$ like all three. How many pupils like none of the three?
\end{probbox}
\probsol{
At least one: $20+18+16-7-6-5+2=38$.\par None: $40-38=2$ pupils.
\dapso{$2$ pupils.}}
\gap
\begin{probbox}{18}{\relax}
In a class of $30$, $15$ like singing, $12$ dancing and $10$ drawing; $5$ like singing and dancing, $4$ dancing and drawing, $3$ singing and drawing, and $1$ likes all three.\par
\pa{a} How many like only singing? \quad \pa{b} How many like none of them?
\end{probbox}
\probsol{
\pa{a} Singing and dancing only: $5-1=4$; singing and drawing only: $3-1=2$.\par Only singing: $15-4-2-1=8$.\par
\pa{b} At least one: $15+12+10-5-4-3+1=26$; none: $30-26=4$.
\dapso{\pa{a} $8$. \pa{b} $4$.}}
\gap
\begin{probbox}{19}{\relax}
In a group, twice as many pupils like football as like volleyball. $6$ pupils like both, and $30$ pupils like at least one of the two games. How many like football and how many like volleyball?
\end{probbox}
\probsol{
Football $+$ volleyball $=30+6=36$.\par Volleyball $1$ part, football $2$ parts: $3$ parts $=36$, one part $=12$.\par
Volleyball: $12$; football: $24$.
\dapso{Football $24$, volleyball $12$.}}
\gap
\begin{probbox}{20}{\relax}
An, Bình and Cường have $60$ marbles. An gives Bình $5$ marbles, Bình gives Cường $8$, then Cường gives An $3$. Now they all have the same number. How many did each have at first?
\end{probbox}
\probsol{
At the end each has $60:3=20$.\par
An: $20-3+5=22$.\par Bình: $20+8-5=23$.\par Cường: $20-8+3=15$.\par Check: $22+23+15=60$.
\dapso{An $22$, Bình $23$, Cường $15$.}}
\gap
\begin{probbox}{21}{\relax}
Two boxes hold $48$ pens. After $7$ pens are moved from box A to box B, the two boxes hold the same number. How many pens were in each box at first?
\end{probbox}
\probsol{
At the end each box holds $48:2=24$.\par Box A: $24+7=31$; box B: $24-7=17$.
\dapso{Box A $31$, box B $17$.}}
\gap
\begin{probbox}{22}{\relax}
Three shelves hold $90$ books. $8$ books are moved from shelf I to shelf II, then $12$ from shelf II to shelf III, then $5$ from shelf III to shelf I. Now the shelves hold the same number. How many books were on each shelf at first?
\end{probbox}
\probsol{
At the end: $90:3=30$ each.\par
Shelf I: $30-5+8=33$.\par Shelf II: $30+12-8=34$.\par Shelf III: $30-12+5=23$.
\dapso{$33$, $34$, $23$ books.}}
\gap
\begin{probbox}{23}{\relax}
Hà gives Lan $4$ stickers, and then Lan gives Hà $10$ stickers. Now Hà has $26$ and Lan has $18$. How many did each have at first?
\end{probbox}
\probsol{
Before Lan gives: Hà $26-10=16$, Lan $18+10=28$.\par
Before Hà gives (at first): Hà $16+4=20$, Lan $28-4=24$.
\dapso{Hà $20$, Lan $24$.}}
\gap
\begin{probbox}{24}{\relax}
From basket A, Mum moves into basket B as many oranges as basket B already has. Now each basket holds $30$ oranges. How many were in each basket at first?
\end{probbox}
\probsol{
Basket B doubled, so at first it held $30:2=15$.\par Basket A gave $15$: at first $30+15=45$.
\dapso{Basket A $45$, basket B $15$.}}
\gap
\begin{probbox}{25}{\relax}
Two tanks hold $80$ litres of water. Tank A pours into tank B as much water as tank B has; then tank B pours into tank A as much as tank A then has. Now the tanks hold the same amount. How much did each tank hold at first?
\end{probbox}
\probsol{
At the end: $40$ litres each.\par
Before step 2: A $40:2=20$, B $40+20=60$.\par
Before step 1: B $60:2=30$, A $20+30=50$.
\dapso{Tank A $50$ litres, tank B $30$ litres.}}
\gap
\begin{probbox}{26}{\relax}
Hùng puts $36$ marbles into three boxes. He moves from box 1 into box 2 as many marbles as box 2 has, then from box 2 into box 3 as many as box 3 has. Now the three boxes hold the same number. How many marbles were in each box at first?
\end{probbox}
\probsol{
At the end: $36:3=12$ each.\par
Before step 2: box 3 had $12:2=6$, box 2 had $12+6=18$, box 1 had $12$.\par
Before step 1: box 2 had $18:2=9$, box 1 had $12+9=21$.
\dapso{$21$, $9$, $6$ marbles.}}
\gap
\begin{probbox}{27}{\relax}
Hòa, Bình and Việt have $70$ cards. Hòa gives Bình $5$ cards, Bình gives Việt $7$, then Việt gives Hòa $9$. Now Hòa and Bình have the same number, and each has twice as many as Việt. How many cards did each have at first?
\end{probbox}
\probsol{
At the end: Việt $1$ part, Hòa $2$ parts, Bình $2$ parts; $5$ parts $=70$, so one part $=14$.\par
End: Hòa $28$, Bình $28$, Việt $14$.\par
At first: Hòa $28-9+5=24$; Bình $28+7-5=30$; Việt $14-7+9=16$.
\dapso{Hòa $24$, Bình $30$, Việt $16$.}}
\gap
\begin{probbox}{28}{\relax}
A woman sells half of her mangoes and $2$ more; then she sells half of the rest and $3$ more. She has $5$ mangoes left. How many mangoes did she have?
\end{probbox}
\probsol{
Before the second sale: $(5+3)\times2=16$ mangoes.\par
At first: $(16+2)\times2=36$ mangoes.
\dapso{$36$ mangoes.}}
\gap
\begin{probbox}{29}{\relax}
Ân and Bảo have $64$ marbles. Ân gives Bảo as many as Bảo has; then Bảo gives Ân as many as Ân has; then Ân again gives Bảo as many as Bảo has. Now they have the same number. How many did each have at first?
\end{probbox}
\probsol{
At the end: $32$ each.\par
Before step 3: Bảo $16$, Ân $48$.\par
Before step 2: Ân $24$, Bảo $40$.\par
At first: Bảo $20$, Ân $44$.
\dapso{Ân $44$, Bảo $20$.}}
\gap
\begin{probbox}{30}{\relax}
On the first day $\dfrac14$ of the water in a tank is used; on the second day $\dfrac13$ of what is left is used. $60$ litres remain. How much water was in the tank at first?
\end{probbox}
\probsol{
$60$ litres is $1-\dfrac13=\dfrac23$ of what was left after day one: $60:\dfrac23=90$ litres.\par
$90$ litres is $1-\dfrac14=\dfrac34$ of the tank: $90:\dfrac34=120$ litres.
\dapso{$120$ litres.}}
\dividerbar{Challenge -- Extension}
\begin{probbox}{31}{\stars{2}}
How many numbers from $1$ to $30$ are divisible by $2$, by $3$ or by $5$? How many are divisible by none of them?
\end{probbox}
\probsol{
By $2$: $15$; by $3$: $10$; by $5$: $6$.\par
By $6$: $5$; by $10$: $3$; by $15$: $2$; by $30$ (all three): $1$.\par
Fill in from the centre:
\begin{center}
\stpVennIII{by $2$}{by $3$}{by $5$}{$8$}{$4$}{$2$}{$4$}{$2$}{$1$}{$1$}{none: $8$}
\end{center}
At least one: $15+10+6-5-3-2+1=22$.\par None: $30-22=8$.
\dapso{$22$ numbers; $8$ numbers.}}
\gap
\begin{probbox}{32}{\stars{1}}
In a class of $40$, $25$ pupils like Maths and $28$ like Literature. What is the least number of pupils who can like both? The greatest?
\end{probbox}
\probsol{
At most $40$ pupils like at least one, so both $\ge25+28-40=13$ (when everybody likes at least one).\par
Both cannot be more than the smaller group: at most $25$ (when every Maths lover also likes Literature).
\dapso{Least $13$, greatest $25$.}}
\gap
\begin{probbox}{33}{\stars{1}}
$45$ pupils each belong to at least one of three clubs A, B, C. Club A has $24$ members, B has $22$, C has $20$; $9$ pupils are in A and B, $7$ in B and C, $8$ in C and A. How many are in all three clubs?
\end{probbox}
\probsol{
$45=24+22+20-9-7-8+(\text{all three})=42+(\text{all three})$.\par
All three: $45-42=3$ pupils.
\dapso{$3$ pupils.}}
\gap
\begin{probbox}{34}{\stars{2}}
In a school contest team, $30$ pupils do Maths, $24$ do English and $20$ do Science. $46$ pupils do Maths or English, $38$ do English or Science, $42$ do Science or Maths, and nobody does all three.\par
\pa{a} How many do Maths and English? \quad \pa{b} English and Science? \quad \pa{c} Science and Maths?\par
\pa{d} How many pupils are in the team?
\end{probbox}
\probsol{
\pa{a} $30+24-46=8$.\par
\pa{b} $24+20-38=6$.\par
\pa{c} $20+30-42=8$.\par
\pa{d} $30+24+20-8-6-8=52$ pupils.
\dapso{\pa{a} $8$. \pa{b} $6$. \pa{c} $8$. \pa{d} $52$ pupils.}}
\gap
\begin{probbox}{35}{\stars{1}}
In a class, $8$ pupils like both dogs and cats, $5$ like neither, and $14$ like dogs. The number who like only cats is twice the number who like only dogs. How many pupils are in the class, and how many like cats?
\end{probbox}
\probsol{
Only dogs: $14-8=6$.\par Only cats: $6\times2=12$.\par
Class: $6+12+8+5=31$ pupils.\par Cats: $12+8=20$ pupils.
\dapso{$31$ pupils; $20$ like cats.}}
\gap
\begin{probbox}{36}{\stars{1}}
How many numbers from $1$ to $100$ are divisible by $3$ but not by $4$?
\end{probbox}
\probsol{
By $3$: $33$ numbers (since $100=3\times33+1$).\par By both $3$ and $4$, i.e.\ by $12$: $8$ numbers ($12,\ 24,\ \dots,\ 96$).\par
By $3$ but not $4$: $33-8=25$.
\dapso{$25$ numbers.}}
\gap
\begin{probbox}{37}{\stars{2}}
Each pupil in a group does at least one of three sports A, B, C: $55$ do A, $45$ do B and $40$ do C. $20$ pupils do exactly two sports and $5$ do all three. How many pupils are in the group?
\end{probbox}
\probsol{
In $55+45+40=140$, a pupil with one sport is counted once, with exactly two sports twice, with three sports three times.\par
Pupils with exactly one sport: $140-2\times20-3\times5=85$.\par
Group: $85+20+5=110$ pupils.
\dapso{$110$ pupils.}}
\gap
\begin{probbox}{38}{\stars{1}}
Of $50$ pupils, $35$ can swim and $30$ can play chess.\par
\pa{a} At least how many pupils can do both? \quad \pa{b} If exactly $4$ pupils can do neither, how many can do both?
\end{probbox}
\probsol{
\pa{a} At most $50$ can do at least one, so both $\ge35+30-50=15$.\par
\pa{b} At least one: $50-4=46$; both: $35+30-46=19$.
\dapso{\pa{a} $15$. \pa{b} $19$.}}
\gap
\begin{probbox}{39}{\stars{2}}
Quân, Sơn and Tùng have $48$ marbles. Quân takes half of his marbles and shares them equally between the other two; then Sơn does the same with half of the marbles he then has; finally Tùng does the same. Now the three have the same number. How many did each have at first?
\end{probbox}
\probsol{
At the end: $16$ each. A giver keeps half; each of the others receives a quarter of what the giver had.
\begin{center}
\renewcommand{\arraystretch}{1.25}
\begin{tabular}{|l|c|c|c|}
\hline
 & Quân & Sơn & Tùng\\\hline
At the end & $16$ & $16$ & $16$\\\hline
Before Tùng gives & $16-8=8$ & $16-8=8$ & $16\times2=32$\\\hline
Before Sơn gives & $8-4=4$ & $8\times2=16$ & $32-4=28$\\\hline
At first & $4\times2=\mathbf{8}$ & $16-2=\mathbf{14}$ & $28-2=\mathbf{26}$\\\hline
\end{tabular}
\end{center}
\dapso{Quân $8$, Sơn $14$, Tùng $26$.}}
\gap
\begin{probbox}{40}{\stars{2}}
There are three drums of oil A, B, C. The oil in drum A is split into $3$ equal parts: one part is poured into B and one into C. Then the oil now in B is split into $4$ equal parts: one part is poured into A and one into C. Then the oil now in C is split into $3$ equal parts: one into A and one into B. Now A holds $40$ litres, B $46$ litres and C $18$ litres. How much oil was in each drum at first?
\end{probbox}
\probsol{
Before C pours: C kept $1$ part of $3$, so C had $18\times3=54$; A and B each got $18$: A $40-18=22$, B $46-18=28$.\par
Before B pours: B kept $2$ parts of $4$, so B had $28\times2=56$; A and C each got $56:4=14$: A $22-14=8$, C $54-14=40$.\par
Before A pours: A kept $1$ part of $3$, so A had $8\times3=24$; B and C each got $8$: B $56-8=48$, C $40-8=32$.
\dapso{A $24$ litres, B $48$ litres, C $32$ litres.}}
\gap
\begin{probbox}{41}{\stars{1}}
Lan reads a book. On day one she reads $\dfrac13$ of the book and $4$ more pages; on day two, half of the remaining pages and $6$ more; on day three, the last $24$ pages. How many pages has the book?
\end{probbox}
\probsol{
On day two, half of the remaining pages and $6$ more leave $24$ pages, so before day two there were $(24+6)\times2=60$ pages left.\par
Day one: $60+4=64$ pages is $1-\dfrac13=\dfrac23$ of the book.\par Book: $64:\dfrac23=96$ pages.
\dapso{$96$ pages.}}
\gap
\begin{probbox}{42}{\stars{2}}
Three boxes I, II, III hold some sweets. First, box I gives boxes II and III each as many sweets as each already holds. Then box II gives boxes I and III each as many as each then holds. Now box I has $20$ sweets, box II $30$ and box III $40$. How many sweets did each box hold at first?
\end{probbox}
\probsol{
Before box II gives: I $20:2=10$, III $40:2=20$, II $30+10+20=60$.\par
Before box I gives: II $60:2=30$, III $20:2=10$, I $10+30+10=50$.\par
Check: $(50,30,10)\to(10,60,20)\to(20,30,40)$.
\dapso{I $50$, II $30$, III $10$ sweets.}}
\gap
\begin{probbox}{43}{\stars{1}}
Bình spends $\dfrac14$ of his money on books, then $\dfrac25$ of the remaining money on a bag, then $15\,000$ dong on snacks. He has $30\,000$ dong left. How much money did he have at first?
\end{probbox}
\probsol{
Before the snacks: $30\,000+15\,000=45\,000$ dong.\par
This is $1-\dfrac25=\dfrac35$ of the money left after the books: $45\,000:\dfrac35=75\,000$ dong.\par
This is $\dfrac34$ of his money: $75\,000:\dfrac34=100\,000$ dong.
\dapso{$100\,000$ dong.}}
\gap
\begin{probbox}{44}{\stars{1}}
Mai says: ``Multiply my age by $3$, add $12$, divide by $5$ and take away $4$: you get $5$.'' How old is Mai?
\end{probbox}
\probsol{
$5+4=9$; \ $9\times5=45$; \ $45-12=33$; \ $33:3=11$.\par
Check: $11\times3+12=45$, $45:5=9$, $9-4=5$.
\dapso{$11$ years old.}}
\gap
\begin{probbox}{45}{\stars{2}}
$40$ pupils each join at least one of three clubs: Reading $22$, Science $18$, Dance $20$. $3$ pupils join all three. How many pupils join exactly two clubs? Exactly one?
\end{probbox}
\probsol{
In $22+18+20=60$: one club counted once, two clubs twice, three clubs three times.\par
Without the $3$ pupils in all three clubs: $37$ pupils, counted $60-9=51$ times.\par
Each pupil in exactly two clubs adds $1$ extra: exactly two $=51-37=14$.\par Exactly one: $37-14=23$.
\dapso{Exactly two: $14$; exactly one: $23$.}}
\gap
\begin{probbox}{46}{\stars{2}}
In a class, $\dfrac35$ of the pupils like Maths, $\dfrac12$ like English, $\dfrac14$ like both, and $6$ pupils like neither. How many pupils are in the class?
\end{probbox}
\probsol{
At least one: $\dfrac35+\dfrac12-\dfrac14=\dfrac{17}{20}$ of the class.\par
Neither: $1-\dfrac{17}{20}=\dfrac{3}{20}$ of the class $=6$ pupils.\par
Class: $6:\dfrac{3}{20}=40$ pupils.
\dapso{$40$ pupils.}}
\gap
\begin{probbox}{47}{\stars{1}}
A shop sells $40\%$ of its rice in the morning and $50\%$ of the rest in the afternoon. $66$ kg are left. How much rice did the shop have?
\end{probbox}
\probsol{
After the morning, $60\%$ is left; the afternoon leaves half of that: $30\%$ of the rice.\par
$30\%$ is $66$ kg, so the rice: $66:30\times100=220$ kg.
\dapso{$220$ kg.}}
\gap
\begin{probbox}{48}{\stars{1}}
An, Bình, Chi and Dũng have $100$ cards. An gives Bình $3$, Bình gives Chi $5$, Chi gives Dũng $7$, and Dũng gives An $9$. Now all four have the same number. How many did each have at first?
\end{probbox}
\probsol{
At the end: $100:4=25$ each.\par
An: $25-9+3=19$; \ Bình: $25+5-3=27$; \ Chi: $25+7-5=27$; \ Dũng: $25-7+9=27$.
\dapso{An $19$, Bình $27$, Chi $27$, Dũng $27$.}}
\gap
\begin{probbox}{49}{\stars{1}}
How many two-digit numbers contain at least one digit $5$?
\end{probbox}
\probsol{
Tens digit $5$: $50$ to $59$, $10$ numbers.\par Ones digit $5$: $15,\ 25,\ \dots,\ 95$, $9$ numbers.\par Both: only $55$.\par
At least one digit $5$: $10+9-1=18$.
\dapso{$18$ numbers.}}
\gap
\begin{probbox}{50}{\stars{2}}
An, Bình and Chi play three games. The loser of each game gives each of the other two as many marbles as that player already has. An loses game 1, Bình loses game 2 and Chi loses game 3. At the end An has $16$, Bình $24$ and Chi $32$ marbles. How many did each have at first?
\end{probbox}
\probsol{
Going back, the two winners halve and the loser gets it back.
\begin{center}
\renewcommand{\arraystretch}{1.25}
\begin{tabular}{|l|c|c|c|}
\hline
 & An & Bình & Chi\\\hline
At the end & $16$ & $24$ & $32$\\\hline
Before game 3 & $8$ & $12$ & $32+8+12=52$\\\hline
Before game 2 & $4$ & $12+4+26=42$ & $26$\\\hline
At first & $4+21+13=\mathbf{38}$ & $\mathbf{21}$ & $\mathbf{13}$\\\hline
\end{tabular}
\end{center}
\dapso{An $38$, Bình $21$, Chi $13$ marbles.}}
